\section{Introduction}
\label{sec:introduction}

Kohn-Sham density functional theory (DFT) is the standard method of electronic-structure simulation because it offers a practical compromise between physical fidelity and computational cost across chemistry, condensed-matter physics, and materials science \cite{HoKo1964,KoSh1965}. Even so, repeated self-consistent solution of the Kohn-Sham equations remains a major bottleneck in workflows that require large supercells, long molecular-dynamics trajectories, broad compositional screening, or repeated re-evaluation under changing thermodynamic conditions. This cost has motivated sustained interest in machine-learned surrogates for atomistic and quantum-mechanical simulation, including direct learning potential energy surfaces, interatomic forces, electron densities, wave-function-related quantities, and reduced electronic representations~\cite{Cangi:2025aa}.
Comment: Add more references: MALA papers and others
Comment: Also cite our ML for DFT review paper.
\cite{BrVo2017,chandrasekaran_solving_2019,TsMi2020,Shao_2023}.

Comment: Extend the paragraph on prior work on MLIPs and add more citations.
Many successful atomistic machine-learning approaches bypass the explicit electronic structure and instead learn scalar or vector observables such as energies and forces.
This strategy is highly effective for interatomic potentials, and recent work has shown that equivariant neural architectures are especially well suited for energy and force prediction because they can encode rotational structure directly in the model rather than relying entirely on handcrafted invariants or data augmentation \cite{KoKoBe22,Rackers_2023}. However, learning only energies and forces also limits what the surrogate can expose: electronic observables, basis-dependent quantities, operator-level post-processing, and Hamiltonian-derived band-structure information are no longer directly available. For applications that remain fundamentally electronic-structure-centered, it is attractive to learn the operators themselves.
Comment: Add examples when the electronic structure is needed for modeling.

\texttt{Mandala} is designed around this operator-learning viewpoint. The framework targets block-sparse matrix quantities arising in localized-orbital electronic-structure calculations, in particular the Hamiltonian, overlap, and density matrices. Instead of treating these objects as dense global arrays, \texttt{Mandala} represents them as atom-pair-resolved sparse blocks and combines this representation with E(3)-equivariant graph neural networks. This yields a learning problem that respects both the locality of the underlying electronic structure and the geometric transformation laws imposed by three-dimensional space.

The software is motivated by three specific design goals. First, it should be modular at the level of both scientific abstractions and code organization. Data ingestion, basis conversion, graph construction, irrep mapping, model definition, and training logic should be independently replaceable while still composing into a single workflow. Second, the framework should support supervision beyond raw matrix reconstruction. When Hamiltonian, overlap, and density matrices are predicted jointly, physically relevant observables such as energies and electron counts can be computed directly from the predicted operators, atomic forces can be obtained through differentiation with respect to atomic coordinates, and for periodic systems the same pathway extends naturally to stress tensors. Third, the framework should make architecture exploration an inherent capability rather than an afterthought. Choices such as irreducible feature content, message-passing depth, tensor-product design, head structure, residual connections, normalization, and observable losses must all be accessible through a coherent configuration and sweep interface.
Comment: Cite deepH papers and other papers that learn the Hamiltonian matrix.

These goals distinguish \texttt{Mandala} from earlier machine-learning workflows in two ways. Relative to purely observable-driven interatomic models, \texttt{Mandala} preserves an operator-level representation that remains useful for electronic-structure analysis. Relative to monolithic research scripts built around one specific dataset or one fixed neural architecture, it provides a reusable framework in which data formats, physical targets, and model variants can be exchanged without rewriting the entire code path. In this sense, \texttt{Mandala} is intended as a unified software environment for sparse operator learning, observable guidance, and architecture exploration.

The following sections describe the sparse matrix data model, the E(3)-equivariant learning formulation, the observable-aware loss construction, the force-training pathway, and the software abstractions that make these components composable. Numerical benchmarks on [TODO: list benchmark systems] and application-specific results for [TODO: list target applications] are reported in Section [TODO: numerical results section].
Comment: Let's decide how to present the results...
